\begin{theorem} \label{thm3}
 Assume that {\rm (A1)--(A3)} hold.
Then, for each $ (u_0,v_0) \in \left( H^2(\Omega) \cap H^1_0(\Omega)\right)^2$ and
$(u_1,v_1) \in \left( H^1_0(\Omega)\right)^2$, problem \eqref{e1.1} has a unique strong solution
$(u, v)$ satisfying, for every $T > 0$,
\begin{gather*}
u, v \in L^{\infty}\left( [0, T), H^2(\Omega) \cap H^1_0(\Omega)\right),\\
u_{t}, v_{t} \in L^2\left( (0, T), H^1_0(\Omega)\right)\cap H^1_0 \left( (0, T)\times \Omega\right),\\
u_{tt}, v_{tt} \in L^{\infty}\left( (0, T), L^2(\Omega)\right).
\end{gather*}
 \end{theorem}

\begin{proof} \quad

\noindent\textbf{Step 1.} Faedo-Galerkin approximation.
We take
\(\{ \phi_j \}_{j=1}^\infty\)
to be the eigenfunctions of the Laplacian operator subject to Dirichlet boundary conditions, satisfying $-\Delta \phi_i = \lambda_i \phi_i $. Then
$\{\phi_j\}_{j=1}^\infty$ forms an orthogonal basis of $ H_0^1(\Omega)$,
\( H^2(\Omega) \cap H_0^1(\Omega) \), and \( L^2(\Omega) \).
Let $V_{N}$ be spanned by the $N$ first functions $ \phi_1, \phi_2, \dots, \phi_N $.
We seek  an approximate solution $(u,v)$ of the form
\[
u^N(x,t) = \sum_{j=1}^N a_j(t) \phi_j(x), \quad
v^N(x,t) = \sum_{j=1}^N b_j(t) \phi_j(x), \quad t\in (0,T),
\]
which satisfies
\begin{gather*}
\int_\Omega u_{tt}^N \phi_j dx + \int_\Omega \nabla u^N .\nabla \phi_j dx - \int_\Omega \int_0^t g(t-s) \nabla u^N(s) .\nabla \phi_j\, ds\,dx+ \alpha \int_\Omega v^N \phi_j \, dx = 0,\\
\int_\Omega v_{tt}^N \phi_j dx + \int_\Omega \nabla v^N .\nabla \phi_j dx+ \alpha \int_\Omega u^N \phi_j \, dx = 0,\quad 1\leq j\leq N
\label{e3.1}
\end{gather*}
and
\begin{equation} \label{e3.2}
\begin{gathered}
 u^{N}(0) = u^{N}_0= \sum_{i=1}^{N} \langle u_0, \phi_i \rangle \phi_i
 \to u_0, \\
  v^{N}(0) = v^{N}_0 = \sum_{i=1}^{N} \langle v_0, \phi_i \rangle \phi_i
 \to v_0 \quad \text{in } H^2(\Omega) \cap H_0^1(\Omega), \\
 u_{t}^N(0) = u_1^N = \sum_{i=1}^{N} \langle u_1, \phi_i \rangle \phi_i \to u_{1}, \\
 v_{t}^N(0) = v_1^N = \sum_{i=1}^{N} \langle v_1, \phi_i \rangle \phi_i \to v_{1} \quad \text{in } H^1_0(\Omega).
 \end{gathered}
\end{equation}
Standard results of ordinary differential equations guarantee the existence of
a unique local solution
$ (u^N, v^N )$ of \eqref{e3.1}-\eqref{e3.2} defined on
$ [0, t_{N}), 0 < t_{N} \leq T $.
Our next step is to obtain a priori estimates that enable us to extend the
solution to the entire interval.
\smallskip

\noindent\textbf{Step 2.} A priori estimate 1.
By multiplying the first equation in \eqref{e3.1} by \(a'_j(t)\) and summing the results over \(j\)
from 1 to \(N\), we obtain, for all \(0 < t \leq t_N\),
\begin{equation}
\begin{aligned}
&\frac{d}{dt} \Big( \frac{1}{2} \| u^N_{t}(t) \|^2 + \frac{1}{2} \| \nabla u^N(t) \|^2 \Big)
+ \alpha \int_{\Omega} v^N(t) \, u^N_{t}(t) \, dx
&- \int_0^t g(t-s) \int_{\Omega} \nabla u^N(s) \nabla u^N_{t}(t) \, dx ds =0.
\end{aligned}\label{e3.3}
\end{equation}
Using Lemma \ref{lem1} for the last term of the left-hand side of \eqref{e3.3}
we find that
\begin{equation}
\begin{aligned}
&\frac{d}{2dt} \Big( \| u^N_{t}(t) \|^2
+ \Big[ 1-\int_0^t g(s) ds \Big] \| \nabla u^N(t) \|^2 + (g\circ \nabla u^N)(t) \Big)
+\alpha \int_{\Omega} v^N(t) \, u^N_{t}(t) dx \\
& +\frac{1}{2} g(t)\parallel \nabla u^N(t) \parallel^2 - \frac{1}{2}( g'\circ\nabla u^N) (t) =0.
\end{aligned} \label{e3.4}
 \end{equation}
Multiplying the second equation in \eqref{e3.1} by $b'_j(t)$ and summing with respect to $j$, we obtain
 \begin{equation}
 \frac{d}{2 dt} \Big( \| v^N_{t}(t) \|^2 + \| \nabla v^N (t) \|^2 \Big)
 + \alpha \int_{\Omega} u^N(t) \, v^N_{t}(t) \, dx =0, \quad 0 < t \leq t_N.
 \label{e3.5}
\end{equation}
Then, combining \eqref{e3.4} and \eqref{e3.5} yields
\begin{equation}
\begin{aligned}
 &\frac{d}{2 dt} \Big( \| u^N_{t}(t) \|^2 + \Big[ 1-\int_0^t g(s)ds\Big]
 \| \nabla u^N(t) \|^2 +(g\circ\nabla u^N)(t)+\| v^N_{t}(t) \|^2 + \| \nabla v^N(t) \|^2 \\
 &+ 2 \alpha \int_{\Omega} v^N(t) u^N(t) dx\Big) \\
&=-\frac{1}{2} g(t)\parallel \nabla u^N(t) \parallel^2 + \frac{1}{2}( g'\circ\nabla u^N )(t).
\end{aligned}\label{e3.6}
\end{equation}
By assumptions (A1)--(A3) we obtain
 \begin{equation}
\begin{aligned}
&\frac{d}{2 dt} \Big( \| u^N_{t} (t) \|^2 + \Big[ 1-\int_0^t g(s)ds\Big]
 \| \nabla u^N(t) \|^2 +(g\circ\nabla u^N)(t)+\| v^N_{t}(t) \|^2
 + \| \nabla v^N (t)\|^2 \\
&+ 2 \alpha \int_{\Omega} v^N(t) u^N(t) dx\Big)\leq 0.
\end{aligned} \label{e3.7}
\end{equation}
Integrating \eqref{e3.7} over $(0, t)$ yields
\begin{equation}
 E_N(t)\leq E_N(0), \label{e3.8}
\end{equation}
where
 \begin{gather*}
 \begin{aligned}
E_N(t)&=\| u^N_{t} (t) \|^2 + \Big[ 1-\int_0^t g(s)ds\Big] \| \nabla u^N(t) \|^2 +(g\circ\nabla u^N)(t)
 +\| v^N_{t}(t) \|^2 + \| \nabla v^N(t) \|^2\\
&\quad + 2 \alpha \int_{\Omega} v^N(t) u^N(t) dx,
\end{aligned}\\
E_N(0)= \| u^N_{1} \|^2 + \| v^N_{1} \|^2 +\| \nabla u^N_0 \|^2
 +\| \nabla v^N_0 \|^2+ 2 \alpha \int_{\Omega} v^N_0 u^N_0 dx.
\end{gather*}
Using Young's and Poincar\'e's inequalities, for $\varepsilon_{1}> 0$, we have
 \begin{equation}
 2\int_{\Omega} v^N(t) u^N(t) dx\geq - \varepsilon_{1} C_p\| \nabla u^N (t) \|^2 - \frac{C_p}{\varepsilon_{1}} \| \nabla v^N (t) \|^2
 \label{e3.9}
 \end{equation}
and
 \begin{equation}
 2 \alpha \int_{\Omega} v^N_0 u^N_0 dx\leq \alpha C_p\| \nabla u^N_0 \|^2 + \alpha C_p \| \nabla v^N_0 \|^2.
 \label{e3.10}
 \end{equation}
Thus
\begin{equation}
 E_N(0)\leq \| u^N_{1} \|^2 + \| v^N_{1} \|^2 +(1+\alpha C_p)\| \nabla u^N_0 \|^2 +(1+\alpha C_p)\| \nabla v^N_0 \|^2.\label{e3.11}
 \end{equation}
Since the sequences $ (u_0^N)_{N \in \mathbb{N}}$,
$ (u_1^N)_{N \in \mathbb{N}}, (v_0^N)_{N \in \mathbb{N}}, (v_1^N)_{N \in \mathbb{N}} $ converge, there exists a positive constant $L_{1}$, independent of $N$, such that
\begin{equation}
E_N(0)\leq L_{1}.\label{e3.12}
\end{equation}
On the other hand,
\begin{equation}
\begin{aligned}
E_N(t)
&\geq \| u^N_{t}(t) \|^2 + \Big[ 1-\int_0^t g(s)ds -\alpha \varepsilon_{1} C_p \Big]
 \| \nabla u^N(t) \|^2 +(g\circ\nabla u^N)(t)+\| v^N_{t} (t) \|^2 \\
 &\quad +\Big( 1-\frac{\alpha C_p}{\varepsilon_{1}}\Big) \| \nabla v^N(t) \|^2.
\end{aligned}\label{e3.13}
\end{equation}
Let us select $\varepsilon_{1}$ such that it satisfies
$\alpha C_p \leq \varepsilon_{1}\leq \frac{l}{\alpha C_p}$.
Then, there exists $C>0$, such that
\begin{equation}
E_N(t)\geq C\left( \| u^N_{t}(t) \|^2 + \| \nabla u^N (t)\|^2
+(g\circ\nabla u^N)(t)+\| v^N_{t}(t) \|^2 + \| \nabla v^N (t)\|^2\right).\label{e3.14}
\end{equation}
From \eqref{e3.8} and \eqref{e3.12}, we deduce
\begin{equation}
\| u^N_{t}(t) \|^2 +\| v^N_{t}(t) \|^2 + \| \nabla u^N(t) \|^2
+ \| \nabla v^N (t)\|^2 +(g\circ\nabla u^N)(t)\leq \frac{L_{1}}{C},\quad
\forall t\leq t_{N}\leq T,
\label{e3.15}
\end{equation}
where $\frac{L_{1}}{C}$ is a constant independent of $t$ and $N$.
\smallskip

\noindent\textbf{A priori estimate 2.}
In \eqref{e3.1}, we substitute $\phi_j $ by $ -\Delta \phi_j$ then
multiply the equation \eqref{e3.1}$_1$ by $a'_j(t)$ and the second equation
\eqref{e3.1}$_2$ by $b'_j(t)$, sum the resulting equations, and use
(A1)--(A3), we arrive at
\begin{align*}
&\frac{d}{2dt}\Big[ \| \nabla u^N_{t}(t) \|^2
+\Big( 1 - \int_0^t g(s)\,ds \Big) \| \Delta u^N(t) \|^2
+ (g \circ \Delta u^N)(t) + \| \nabla v^N_{t}(t) \|^2
+ \| \Delta v^N(t) \|^2 \\
& +2 \alpha \int_{\Omega} \nabla v^N(t) \nabla u^N(t) dx\Big]
-\frac{1}{2} \left( g' \circ \Delta u^N \right)(t)
+\frac{1}{2} g(t) \| \Delta u^N(t) \|^2\leq 0.
\end{align*}

Applying Young's and Poincar\'e's inequalities, we have for $\varepsilon_{3}> 0$,
 \begin{equation*}
 2\int_{\Omega} \nabla v^N(t) \nabla u^N(t) dx
 \geq - \varepsilon_{3} C_p\| \Delta u^N(t) \|^2 - \frac{C_p}{\varepsilon_{3}}\| \Delta v^N(t) \|^2,
 \end{equation*}
and
\[
 2 \int_{\Omega} \nabla v^N_0 \nabla u^N_0 dx
 \leq C_p\| \Delta u^N_0 \|^2 +C_p \| \Delta v^N_0 \|^2.
\]
Then
\begin{equation}
\begin{aligned}
& \| \nabla u^N_{1} \|^2+ \| \nabla v^N_{1} \|^2+\| \Delta u^N_0 \|^2
 + \| \Delta v^N_0 \|^2 +2 \alpha \int_{\Omega} \nabla v^N_0 \nabla u^N_0 dx \\
&\leq \| \nabla u^N_{1} \|^2+ \| \nabla v^N_{1} \|^2
 +(1+\alpha C_p) \| \Delta u^N_0 \|^2
 +(1+\alpha C_p)\| \Delta v^N_0 \|^2.
\end{aligned} \label{e3.16}
\end{equation}
Since the sequences $ (u_0^N)_{N \in \mathbb{N}}$,
$ (u_1^N)_{N \in \mathbb{N}}, (v_0^N)_{N \in \mathbb{N}}, (v_1^N)_{N \in \mathbb{N}} $
converge, there exists a positive constant $L_2$, independent of $N$, such that
\begin{equation}
 \| \nabla u^N_{1} \|^2+ \| \nabla v^N_{1} \|^2
 +\| \Delta u^N_0 \|^2 + \| \Delta v^N_0 \|^2
 +2 \alpha \int_{\Omega} \nabla v^N_0 \nabla u^N_0 dx\leq L_2. \label{e3.17}
\end{equation}
Furthermore,
\begin{align*}
&\| \nabla u^N_{t}(t) \|^2+\Big( 1 - \int_0^t g(s)\,ds \Big) \| \Delta u^N(t) \|^2
 + (g \circ \Delta u^N)(t) + \| \nabla v^N_{t}(t) \|^2+ \| \Delta v^N (t)\|^2 \\
&+2 \alpha \int_{\Omega} \nabla v^N (t)\nabla u^N(t) dx \\
&\geq \| \nabla u^N_{t}(t) \|^2+\Big( 1 - \int_0^t g(s)\,ds
 - \varepsilon_{3} C_p \Big) \| \Delta u^N(t) \|^2 \\
&\quad + (g \circ \Delta u^N)(t) + \| \nabla v^N_{t}(t) \|^2
 +\Big( 1-\frac{\alpha C_p}{\varepsilon_{3}}\Big) \| \Delta v^N(t) \|^2.
\end{align*}

We choose $\varepsilon_{3}$ such that it satisfies
$\alpha C_p \leq \varepsilon_{3}\leq \frac{l}{\alpha C_p}$.
Then there exists $C>0$, such that
\begin{equation}
\begin{aligned}
&\| \nabla u^N_{t}(t) \|^2+\Big( 1 - \int_0^t g(s)\,ds \Big) \| \Delta u^N(t) \|^2
+ (g \circ \Delta u^N)(t) + \| \nabla v^N_{t}(t) \|^2+ \| \Delta v^N(t) \|^2\\
&+2 \alpha \int_{\Omega} \nabla v^N (t)\nabla u^N (t) dx \\
&\geq C \Big( \| \nabla u^N_{t}(t) \|^2+ \| \Delta u^N (t)\|^2
+ (g \circ \Delta u^N)(t) + \| \nabla v^N_{t}(t) \|^2
+ \| \Delta v^N (t)\|^2 \Big).
\end{aligned} \label{e3.18}
\end{equation}
From \eqref{e3.16}, \eqref{e3.17} and \eqref{e3.18}, it follows that
\begin{equation}
\| \nabla u^N_{t}(t) \|^2+ \| \Delta u^N(t) \|^2 + (g \circ \Delta u^N)(t)
+ \| \nabla v^N_{t}(t) \|^2+ \| \Delta v^N(t) \|^2
\leq \frac{L_2}{C}, \; \forall\, t\leq t_{N} \leq T,\label{e3.19}
 \end{equation}
where $\frac{L_2}{C}$ is a constant independent of $t$ and $N$.
\smallskip

\noindent\textbf{A priori estimate 3.}
 By multiplying the first equation in \eqref{e3.1}$_1$ by
 \(a''_j(t)\) and summing the results over \(j\) from 1 to \(k\), we obtain
\begin{align*}
&\int_\Omega |u_{tt}^N(t)|^2 dx \\
& = \int_\Omega \Delta u^N(t) u_{tt}^N(t) \, dx
 - \int_\Omega \int_0^t g(t-s) \Delta u^N(s) u_{tt}^N(t) \, ds\,dx
 - \alpha \int_\Omega v^N(t) u_{tt}^N(t) \, dx\\
& = \int_\Omega \Delta u^N(t) u_{tt}^N(t) \, dx
 + \int_\Omega u_{tt}^N(t) \int_0^t g(t-s) \Big(\Delta u^N(t)- \Delta u^N(s)\Big)\, ds\,dx\\
&\quad -\int_0^t g(s) ds \int_\Omega u_{tt}^N(t) \Delta u^N(t) dx
 - \alpha \int_\Omega v^N(t) u_{tt}^N(t) \, dx\\
&\leq \varepsilon \int_\Omega |u_{tt}^N(t)|^2 dx
 +\frac{1}{4\varepsilon}\int_\Omega |\Delta u^N(t)|^2 \, dx
 + \varepsilon \int_\Omega |u_{tt}^N(t)|^2 dx+\frac{1-l}{4\varepsilon}
 (g \circ \Delta u^N)(t)\\
&\quad +\varepsilon(1-l) \int_\Omega |u_{tt}^N(t)|^2 dx
 + \frac{1-l}{4\varepsilon} \int_\Omega |\Delta u^N(t)|^2dx
 + \varepsilon \int_\Omega |u_{tt}^N(t)|^2 dx
 +\frac{\alpha^2}{4\varepsilon}\int_\Omega |v^N(t)|^2 dx \\
&\leq \varepsilon (4-l)\int_\Omega |u_{tt}^N(t)|^2 dx
 +\frac{2-l}{4\varepsilon} \int_\Omega |\Delta u^N(t)|^2dx+
 \frac{1-l}{4\varepsilon} (g \circ \Delta u^N)(t)
 +\frac{\alpha^2}{4\varepsilon}\int_\Omega |v^N(t)|^2 dx.
\end{align*}
Hence, for some $c >0$, we have
\begin{align}
\left(1-\varepsilon(4-l)\right) \int_\Omega |u_{tt}^N(t)|^2 dx
\leq \frac{c}{4\varepsilon}(L_1+L_2).
\end{align}
We take $\varepsilon >0$ small enough, so
\begin{align}
\|u_{tt}^N(t)\|^2 \leq \tilde{c}(L_1+L_2)=L_3. \label{e3.20}
\end{align}
Similarly, we obtain $\|v_{tt}^N(t)\|^2 \leq L_4$, where $L_{3}, L_{4} $
are constants independent of $t$ and $N$.

From \eqref{e3.15}, \eqref{e3.19} and \eqref{e3.20}, we conclude that
\begin{equation} \label{e3.21}
\begin{gathered}
(u^{N})_{N} \text{ and } (v^N)_{N} \text{ are bounded in }
L^\infty((0,T), H^2(\Omega) \cap H_0^1 (\Omega)),\\
(u^N_{t})_{N} \text{ and } (v^N_{t})_{N} \text{ are bounded in }
L^\infty((0,T), H^1_0(\Omega)), \\
(u^N_{tt})_{N} \text{ and } (v^N_{tt})_{N} \text{ are bounded in }
 L^\infty((0,T), L^2(\Omega)).
\end{gathered}
\end{equation}

\noindent\textbf{Step 3.}
From \eqref{e3.21}, there exist subsequences of $(u^{N})_{N}$ and $(v^{N})_{N}$
(still denoted by $(u^{N})_{N}$ and $(v^{N})_{N})$, and two functions
$u, v : \Omega \times [0, T) \to \mathbb{R}$, such that
\begin{equation}
\begin{gathered}
u^{N} \overset{\star}{\rightharpoonup} u \text{ and }
 v^{N} \overset{\star}{\rightharpoonup} v \quad \text{in }
 L^\infty((0,T), H^2(\Omega) \cap H_0^1 (\Omega)), \\
u_{t}^{N} \overset{\star}{\rightharpoonup} u_{t} \text{ and }
 v_{t}^{N} \overset{\star}{\rightharpoonup} v_{t} \text{ in } L^\infty((0,T),
 H^1_0(\Omega)), \\
 u_{tt}^{N} \overset{\star}{\rightharpoonup} u_{tt} \text{ and }
 v_{tt}^{N} \overset{\star}{\rightharpoonup} v_{tt} \text{ in }
 L^\infty((0,T), L^2(\Omega)).
\end{gathered}\label{e3.22}
\end{equation}
Thanks to convergences \eqref{e3.2} and \eqref{e3.22}, taking $N \to +\infty$
in \eqref{e3.1}, we obtain
\begin{equation}
\begin{gathered}
\int_\Omega u_{tt}\phi_j dx + \int_\Omega \nabla u .\nabla \phi_j dx
- \int_\Omega \int_0^t g(t-s) \nabla u(s) .\nabla \phi_j\, ds\, dx
+ \alpha \int_\Omega v \phi_j \, dx = 0,\\
\int_\Omega v_{tt} \phi_j dx + \int_\Omega \nabla v.\nabla \phi_j dx
+ \alpha \int_\Omega u \phi_j \, dx = 0,\quad 1\leq j\leq N.
\end{gathered} \label{e3.23}
\end{equation}
 Consequently, for every $ \phi \in H_0^1(\Omega)\cap H^2(\Omega)$,
\begin{gather*}
\int_\Omega \Big[ u_{tt}- \Delta u + \int_0^t g(t-s) \Delta u(s)\, ds\,dx
 + \alpha v\Big] \phi \, dx = 0,\\
\int_\Omega [ v_{tt} - \Delta v + \alpha u] \phi \, dx = 0.
\end{gather*}
Therefore, $(u,v)$ solves the equations in \eqref{e1.1}.
\smallskip

\noindent\textbf{Step 4. Initial conditions.}
First, by Lions' lemma \cite{Lions} and since
\begin{gather*}
u^{N} \overset{\star}{\rightharpoonup} u \quad \text{in }
 L^\infty((0,T), H^2(\Omega) \cap H_0^1 (\Omega)), \\
u_{t}^{N} \overset{\star}{\rightharpoonup} u_{t} \quad \text{in }
 L^\infty((0,T), H^1_0(\Omega)),
\end{gather*}
we deduce, up to a subsequence, that
\begin{equation}
u^{N} \to u \quad \text{in } C([0, T], H^1_0(\Omega)).\label{e3.24}
\end{equation}
Therefore, \( u^{N}(\cdot, 0) \) is defined and
\[
u^{N}(\cdot, 0) \to u(\cdot, 0) \quad \text{in } H^1_0(\Omega).
\]
But,
\[ u^{N}(\cdot, 0) = u^{N}_0 \to u_0 \,\, \text{in}\,\, H^2(\Omega) \cap H^1_0(\Omega).\] Then \( u(\cdot, 0) = u_0 \, \). Similarly, we obtain $v(\cdot, 0)=v_0. $

Noting that
\[
 u_{tt}^{N} \overset{\star}{\rightharpoonup} u_{tt} \quad \text{and} \quad
 v_{tt}^{N} \overset{\star}{\rightharpoonup} v_{tt} \quad \text{in }
 L^\infty((0,T), L^2(\Omega)),
\]
we obtain $ u_t , v_{t} \in C([0, T), L^2(\Omega))$, as in lions \cite{Lions},
and thus \( u_t(\cdot, 0) \) has a meaning. Furthermore
\[
u^N_t(\cdot, 0) \to u_t(\cdot, 0) \quad \text{and} \quad
v^N_t(\cdot, 0) \to v_t(\cdot, 0) \quad\text{in } L^2(\Omega).
\]
But,
\[ u^N_t(\cdot, 0) = u^N_1 \to u_1 \quad \text{and} \quad
 v^N_t(\cdot, 0) = v^N_1 \to v_1 \quad \text{in } H^1_0(\Omega).
\]
Hence,
$u_t(x, 0) = u_1(x)$ and $v_t(x, 0) = v_1(x)$.
\smallskip

\noindent\textbf{Uniqueness.}
Assume that problem \eqref{e1.1} has two weak solutions, \( (u, v) \) and
\( (\overline{u},\overline{v}) \). Then,
\( (\mathbf{u}, \mathbf{v}) = (u - \overline{u}, v - \overline{v}) \) satisfies
\begin{gather}
 \mathbf{u}_{tt} - \Delta \mathbf{u} + \int_0^t g(t-s) \Delta \mathbf{u} \, ds
 + \alpha \mathbf{v} = 0, \quad \text{in } \Omega \times (0, T), \label{e3.25}\\
 \mathbf{v}_{tt} - \Delta \mathbf{v} + \alpha \mathbf{u} = 0, \quad
 \text{in } \Omega \times (0, T), \label{e3.26} \\
 \mathbf{u}(x,t) = \mathbf{v}(x,t) = 0, \quad \text{on } \partial\Omega \times (0, T),
 \label{e3.27}\\
 \mathbf{u}(x,0) = \mathbf{u}_t(x,0) =0,\, \mathbf{v}(x,0)=\mathbf{v}_t(x,0) = 0,
 \quad \text{in } \Omega. \label{e3.28}
\end{gather}
Multiplying \eqref{e3.25} by $\mathbf{u}_{t}$ and \eqref{e3.26} by $\mathbf{v}_{t}$,
integrating over $\Omega$, and adding the two resulting equations, we find
\begin{equation}
\begin{aligned}
&\frac{1}{2}\frac{d}{dt} \Big( \| \mathbf{u}_{t}(t) \| ^2
 + \Big[ 1-\int_0^t g(s) ds \Big] \| \nabla \mathbf{u} (t)\|^2
 +(g\circ \nabla \mathbf{u})(t) \\
&+2 \alpha \int_\Omega \mathbf{v}\textbf{ u}\,dx
 +\| \mathbf{v}_{t}(t) \|^2 + \| \nabla \mathbf{v} (t)\|^2 \Big)\\
& = -\frac{1}{2} g(t)\parallel \nabla \mathbf{u}(t) \parallel^2
 + \frac{1}{2}( g'\circ\nabla \mathbf{u}) (t).
\end{aligned} \label{e3.29}
\end{equation}
Using assumptions (A1)--(A3), and integrating over $(0, t),\, t\leq T$, we obtain
\[
 \| \mathbf{u}_{t}(t) \| ^2+ \Big[ 1-\int_0^t g(s) ds \Big]
 \| \nabla \mathbf{u}(t) \|^2 +(g\circ \nabla \mathbf{u})(t)
 +2 \alpha \int_\Omega \mathbf{v}\textbf{ u}dx
 +\| \mathbf{v}_{t} (t)\|^2 + \| \nabla \mathbf{v}(t) \|^2 \leq 0.
\]
 From Young and Poincar\'e's inequalities, we obtain
\begin{align*}
&\| \mathbf{u}_{t}(t) \| ^2+ \Big[ 1-\int_0^t g(s) ds
 -\alpha \varepsilon_{1} C_p \Big] \| \nabla \mathbf{u}(t) \|^2
 +(g\circ \nabla \mathbf{u})(t)\\
&+\| \mathbf{v}_{t}(t) \|^2 + \Big( 1-\frac{\alpha C_{p}}{\varepsilon_{1} }\Big)
\| \nabla \mathbf{v}(t) \|^2 \leq 0.
\end{align*}
We select $\varepsilon_{1}$ such that it satisfies\,
$\alpha C_p \leq \varepsilon_{1}\leq \frac{l}{\alpha C_p}$. Then
\[
 \| \mathbf{u}_{t} \|^2 + \| \nabla \mathbf{u} \|^2
 +(g\circ\nabla \mathbf{u})(t)+\| \mathbf{v}_{t} \|^2 + \| \nabla \mathbf{v} \|^2=0.
\]
Hence, \( \mathbf{u}_t(\cdot, t) = \mathbf{v}_t(\cdot, t) = 0 \) and
\( \nabla \mathbf{u}(\cdot, t) = \nabla\mathbf{v}(\cdot, t) = 0 \), for all
\( t \in (0, T) \). This implies that \( \textbf{ u} = \textbf{ v} = 0 \) on
\( \Omega \times (0, T) \), given that \( \mathbf{u} = \mathbf{v} = 0 \) on
\( \partial\Omega \times (0, T) \). This establishes the uniqueness.
\end{proof}

